% Standalone section fragment for the parent article.
% Requires amsmath, amssymb, amsthm; theorem/lemma/proposition/corollary/remark.
% Bibliography keys and complete entries are supplied at the end as comments.
% All labels use the mon: prefix.

\section{Exact arithmetic recognition and local obstructions from Montesinos tangles}
\label{sec:mon:montesinos}

This section provides a complete arithmetic decision procedure on an explicitly
specified family, and a certificate that can reject an arbitrary diagram
containing a suitable local tangle.  The topology underlying the complete
Montesinos classification is classical.  The local unknottability criterion is
due to Nogueira and Salgueiro \cite[Theorem~4.9]{nogueira-salgueiro}.
Our purpose is to state the input model, prove the arithmetic and diagram
verification obligations, and make these results available to the recognizer.
No statement here asserts efficient discovery of a Montesinos presentation for
an arbitrary knot.

\subsection{Conventions and exact component counts}

Write $R(p/q)$ for the rational tangle of fraction $p/q$, with boundary points
$NW,NE,SE,SW$ in their cyclic order.  Horizontal addition joins the east pair
of one tangle to the west pair of the next.  The numerator closure $N(T)$ joins
$NW$ to $NE$ and $SW$ to $SE$ outside the tangle disk.  We use the standard
fraction rules $R(x)+[a]=R(x+a)$ and $R(x)^{-1}=R(1/x)$; the latter operation
is a quarter rotation followed by crossing reversal.  A rotation alone sends
$x$ to $-1/x$.  These conventions, including the distinction between rotation
and inversion, are those of \cite[Section~2]{kauffman-lambropoulou-tangles}.
An implementation must calibrate its crossing convention against these rules.

The structured input is
\[
 M(e;p_1/q_1,\ldots,p_s/q_s)
 =N\bigl([e]+R(p_1/q_1)+\cdots+R(p_s/q_s)\bigr),
\]
where $e$ and the rational-tangle continued-fraction coefficients are signed
binary integers.  Every final slope is finite.  Each pair is reduced, with
$q_i>0$.  A denominator-one summand is absorbed into $e$.  We thereafter write
the remaining slopes as $\beta_i/\alpha_i$, $1\le i\le r$, where
$\alpha_i\ge2$ and $\gcd(\alpha_i,\beta_i)=1$.
It is optional to impose $0<\beta_i<\alpha_i$: replacing
$\beta_i$ by $\beta_i-a\alpha_i$ and replacing $e$ by $e+a$ preserves the
Montesinos link.  This is the standard integer-transfer isotopy; it follows
by moving integral twists through rational summands using the rational flip
isotopies \cite[Section~2]{kauffman-lambropoulou-tangles}.
The tests below do not require the optional remainder normalization. The delivered implementation always performs it.

For a finite continued fraction $[a_1,\ldots,a_\ell]$, the product
\[
 \begin{pmatrix}a_1&1\\1&0\end{pmatrix}\cdots
 \begin{pmatrix}a_\ell&1\\1&0\end{pmatrix}
 \binom10=\binom pq
\]
computes its projective fraction without floating-point arithmetic.  Intermediate
infinite values are harmless in this formulation.  A final $q=0$ is outside
the stated input model and must produce an unsupported-input response, not a
knot verdict.  In particular, one must not replace an infinite tangle by an
integer summand or apply the formulas below to it.

\begin{lemma}[Endpoint parity and component count]\label{lem:mon:parity}
Let $k$ be the number of even denominators among the reduced finite slopes in
a horizontal Montesinos sum.  The sum has $\max(k-1,0)$ closed strand
components in its interior.  Consequently it is a genuine two-string tangle
with no closed components if and only if $k\le1$.
Its numerator closure has
\[
 c=\begin{cases}
 k,&k>0,\\
 1,&k=0\text{ and }e+\sum_i\beta_i\text{ is odd},\\
 2,&k=0\text{ and }e+\sum_i\beta_i\text{ is even}.
 \end{cases}
\]
The same formula includes $r=0$.
\end{lemma}

\begin{proof}
The endpoint pairing of a rational tangle depends only on the coprime pair
$(p,q)$ modulo two.  There are three possibilities:
\[
\begin{array}{c|c|c}
(p\bmod2,q\bmod2)&\text{pairing}&\text{type}\\\hline
(0,1)&NW\!\leftrightarrow\!NE,\ SW\!\leftrightarrow\!SE&H\\
(1,0)&NW\!\leftrightarrow\!SW,\ NE\!\leftrightarrow\!SE&V\\
(1,1)&NW\!\leftrightarrow\!SE,\ SW\!\leftrightarrow\!NE&X.
\end{array}
\]
This can be verified inductively from $[0]$, since adding one integral twist
and inverting the tangle act on the three pairings exactly as they act on
$(p,q)$ modulo two.  Thus it holds for every continued-fraction construction,
including constructions with negative and zero coefficients.

An $H$ or $X$ piece carries the left endpoint pair to the right endpoint pair
by the identity or the transposition.  If there is no $V$ piece, the total
transposition parity is $e+\sum_i\beta_i$; numerator closure gives the cycles
of this two-element permutation.  If there are $k>0$ pieces of type $V$,
each interval between consecutive such pieces closes one component.  There
are $k-1$ such intervals in the open sum.  The two remaining outer caps
close one additional component under numerator closure, giving $k$ components.
No crossings are removed or changed in this argument: it follows the strands
through the actual tangle connectivity.
\end{proof}

The component check is mandatory.  A link with several components is outside
the ordinary knot-recognition domain; its presence is not evidence for a
nontrivial knot.  In a local witness the corresponding mandatory check is that
the selected portion consists of two open strands and no internal closed
strand.  Following opposite darts through the crossings verifies this directly,
independently of the arithmetic parity computation.

\subsection{A direct nontriviality certificate for the branched cover}

The branched double cover of a Montesinos link is Seifert fibered, with a
presentation, in the above sign convention,
\begin{equation}\label{eq:mon:sfs-group}
 \pi_1\Sigma_2(M)=
 \left\langle h,x_1,\ldots,x_r\ \middle|\
 [h,x_i]=1,\ x_i^{\alpha_i}h^{\beta_i}=1,
 \ x_1\cdots x_r=h^e\right\rangle.
\end{equation}
This is the classical Montesinos construction; see the application and
attribution in \cite[proof of Theorem~4.9]{nogueira-salgueiro}.
Here a lift of each rational-tangle ball is a solid torus.  The central piece
is the circle bundle over the punctured sphere, and the meridian of the
$i$th attached solid torus is $\alpha_i x_i+\beta_i h$.  These gluing slopes
give the displayed relations by van Kampen's theorem.

We give an explicit elementary representation to avoid treating the
nontriviality of a triangle group as a group-decision black box.

\begin{lemma}[An explicit nontrivial triangle representation]
\label{lem:mon:triangle}
For all integers $a,b,c\ge2$, the group
\[
 \langle x,y,z\mid x^a=y^b=z^c=xyz=1\rangle
\]
has a nontrivial representation in $\operatorname{PSL}_2(\mathbb C)$.
\end{lemma}

\begin{proof}
Set
\[
 \lambda=e^{\pi i/a},\qquad t=2\cos(\pi/b),\qquad
 s=2\cos(\pi/c),\qquad d=\lambda-\lambda^{-1},
\]
and define
\[
 u=\frac{s-\lambda^{-1}t}{d},\qquad
 v=\frac{\lambda t-s}{d},\qquad
 X=\begin{pmatrix}\lambda&0\\0&\lambda^{-1}\end{pmatrix},\qquad
 Y=\begin{pmatrix}u&1\\uv-1&v\end{pmatrix}.
\]
The denominator $d$ is nonzero because $a\ge2$.  Direct calculation gives
$\det X=\det Y=1$, $\operatorname{tr}Y=t$, and
$\operatorname{tr}(XY)=s$.  Thus the distinct eigenvalues of $Y$ are
$e^{\pi i/b}$ and $e^{-\pi i/b}$, while those of $XY$ are
$e^{\pi i/c}$ and $e^{-\pi i/c}$.  It follows that
\[
 X^a=Y^b=(XY)^c=-I.
\]
Sending $x,y,z$ to the projective classes of $X,Y,(XY)^{-1}$ therefore
satisfies every relation.  The image of $x$ has order exactly $a$, so it is
nontrivial.
\end{proof}

This is a uniform symbolic certificate formula, rather than a requirement to
expand cyclotomic fields of possibly enormous degree.  Its instantiated
parameters are just the three binary integers $a,b,c$.  The verifier checks
that they are the selected denominators and that they are at least two; the
lemma supplies the representation for every such triple.  No numerical cosine
evaluation or unproved finite-group search is involved.

\begin{proposition}[Three genuine exceptional fibers obstruct the unknot]
\label{prop:mon:three}
If $r\ge3$, a Montesinos link of the stated form cannot be the unknot.
\end{proposition}

\begin{proof}
In \eqref{eq:mon:sfs-group}, quotient first by $h=1$, then by $x_i=1$ for
all but any three selected indices, retaining their cyclic order.  The resulting
quotient is the triangle group with orders equal to their denominators.
It is nontrivial by Lemma~\ref{lem:mon:triangle}.  Hence the branched double
cover is not simply connected.  The double cover branched over the unknot
is $S^3$, whose fundamental group is trivial.
\end{proof}

\subsection{A complete decision procedure on the structured input}

Put
\begin{equation}\label{eq:mon:delta}
 Q=\prod_{i=1}^r\alpha_i,\qquad
 D=eQ+\sum_{i=1}^r\beta_i\frac{Q}{\alpha_i},
\end{equation}
with the empty-product convention $Q=1$.  Abelianizing
\eqref{eq:mon:sfs-group} gives a square integer relation matrix whose
determinant has absolute value $|D|$.  Indeed, eliminate the diagonal block
$\operatorname{diag}(\alpha_1,\ldots,\alpha_r)$ over $\mathbb Q$; its Schur
complement is $-e-\sum_i\beta_i/\alpha_i$.  When $D\ne0$, the order of first
homology is $|D|$; when $D=0$, first homology is infinite.

\begin{theorem}[Complete Montesinos recognition in compressed arithmetic]
\label{thm:mon:complete}
For a structured finite-slope Montesinos input, first compute the component
count in Lemma~\ref{lem:mon:parity}.  If it is one, the represented knot is
the unknot exactly when $r\le2$ and $|D|=1$.
\end{theorem}

\begin{proof}
The case $r\ge3$ is Proposition~\ref{prop:mon:three}.
For $r=0$ or $r=1$, the link is the numerator of a rational tangle.
For $r=2$, absorb the integral summand into one rational tangle; the numerator
closure of the sum of two rational tangles is a two-bridge link.  This last
statement concerns the closure, and does not assert that the open sum is
rational; see \cite[Section~2]{kauffman-lambropoulou-tangles}.
The branched cover of a two-bridge link is a lens space, and its numerator
parameter has absolute value $|D|$.  The rational-link classification says
that a two-bridge knot is the unknot precisely when the absolute numerator
is one \cite{kauffman-lambropoulou-knots}.  This proves the claimed equivalence.
Equivalently, the $r=2$ assertion is the integral-closure calculation in
\cite[Theorem~4.9]{nogueira-salgueiro}.
\end{proof}

\begin{proposition}[Bit complexity of the arithmetic implementation]
\label{prop:mon:complexity}
Let $B$ be the complete binary length of a structured description, including
coefficients, signs, and list delimiters. The production arithmetic classifier
uses $O(B^2)$ schoolbook bit operations and $O(B)$ bits of integer data.
Its valid branch certificate has $O(B)$ bits. The independent verifier in the
package has the conservative bound $O(B^3)$ on these certificates. Input
parsing must additionally charge the actual size of a purported certificate.
These are bounds in compressed source length; they do not include PD expansion.
\end{proposition}

\begin{proof}
Give a coefficient $a$ the weight
$b(a)=1+\lceil\log_2(1+|a|)\rceil$ and charge at least one bit per delimiter.
The sum of all weights is $O(B)$. The column recurrence starts at $(p,q)=(1,0)$
and, reading a continued fraction backwards, replaces it by
$(ap+q,p)$. Every matrix has determinant $-1$, so the final pair is primitive
without a gcd computation. The one-norm of a column is bounded by the
product of $(1+|a|)$ over its coefficients; allowing a constant for the
initial state gives $O(B)$ bits for every entry, even with zeros and signs.
The total size of the final numerator--denominator pairs over all input
summands is $O(B)$ as well.

More precisely, a recurrence with current entry length $s$ and coefficient
weight $b$ costs $O(sb+s)$ with schoolbook arithmetic. Summing these costs
is $O(B^2)$, since all partial lengths are bounded by the corresponding
partial sum of coefficient weights and there are $O(B)$ entries.
Changing the sign to make $q>0$ is linear in its size. For each finite pair,
Euclidean division $p=aq+\beta$ computes $0\le\beta<q$. The sum of the
squares of the pair lengths is at most the square of their total length,
so all these divisions together still cost $O(B^2)$. The implementation
absorbs every integral part into $e$ and drops exactly the pairs with $q=1$.

To evaluate \eqref{eq:mon:delta}, begin with $(d,Q)=(e,1)$ and, for each
remaining $(\beta,\alpha)$, update
\[
 (d,Q)\longmapsto(d\alpha+\beta Q,Q\alpha).
\]
No rational cancellation is performed. The total bit lengths of the
$(\beta,\alpha)$ pairs are $O(B)$, and every accumulated entry has
$O(B)$ bits. Weighting each multiplication by both operand lengths shows
that the sum of multiplication costs is $O(B^2)$; the at most $O(B)$
additions and parity checks contribute the same bound. All stored slopes,
accumulated integers, counts, and the certificate together use $O(B)$ bits
of integer and encoded source data. This is an arithmetic-space statement,
not a measurement of Python object overhead or process memory.

The production code tracks the maximum bit length of its recorded
continued-fraction and accumulator states with a constant number of integer
operations per state, costing $O(B^2)$ bit operations in total.
The independent verifier reconstructs each fraction by a matrix product and
recomputes product--sum expressions for the certificate. Its current replay
of the prefix statistics contains nested loops. Bounding their work by
$O(B^3)$ is safe: final products have $O(B)$ bits and the sums of operand
lengths are $O(B)$. It never needs to expand a cyclotomic field, factor an
integer, or enumerate a finite quotient. For valid certificates all compared
fields have total length $O(B)$.
\end{proof}

The statistic \code{max_entry_bits} means the maximum \emph{tracked}
continued-fraction or accumulator entry length. It is not the maximum length
of every multiplication temporary, and is not peak memory. Temporary
products still have $O(B)$ bits by the same estimates.
For instance, the calculation for $M(-2;7/9)$ can form the temporary $-18$
while the tracked final numerator is $-11$; the former has one more binary
digit. The theorem does not infer a space bound from this empirical statistic.

The one-component check can be implemented especially simply as $D$ odd.
If exactly one denominator is even, its term $\beta_i Q/\alpha_i$ is odd
and every other term in \eqref{eq:mon:delta} is even. If at least two are
even, every term is even. If none is even, the parity is
$e+\sum\beta_i$. Lemma~\ref{lem:mon:parity} therefore shows that $D$ is odd
if and only if the numerator closure has one component. This equivalence
includes zero exceptional fibers, and justifies the code's exact component
rejection without expanding a diagram.

\subsection{An Alexander-trivial family of genus-one knots}

For odd integers $p,q,r$, use the conventional pretzel diagram
$P(p,q,r)=M(0;1/p,1/q,1/r)$, with signs of the twist boxes as displayed in
\cite[Figure~1 and Theorem~1(3)]{belousov-pretzel}.  Its usual two-disk,
three-band spanning surface is orientable: all three bands have odd twisting,
so the disk orientations can be chosen compatibly.  It has Euler characteristic
$2-3=-1$ and one boundary component, hence genus one.
Take one oriented cycle through the first and second bands and another through
the second and third bands, choosing their intersection to be $+1$.
In this basis a Seifert matrix is
\begin{equation}\label{eq:mon:pretzel-seifert}
 V=\frac12\begin{pmatrix}p+q&q+1\\q-1&q+r\end{pmatrix}.
\end{equation}
For completeness, the self-pairings are half the signed twists in the two
bands traversed by the corresponding cycle.  In the shared middle band the
off-diagonal pairings are $(q+1)/2$ and $(q-1)/2$: the parallel push-off on
the two sides of the single transverse disk intersection differs by one.
Their difference is the chosen intersection number, as independently required
by $V-V^{\mathsf T}=\left(\begin{smallmatrix}0&1\\-1&0\end{smallmatrix}\right)$.
Reversing the entire diagram convention replaces the corresponding form by
its mirror form and does not change the conclusions below.  There is also
an independent check from the skein-derived formula in
\cite[Theorem~1(3)]{belousov-pretzel}.

Writing
\[
 A=\det V=\frac{pq+pr+qr+1}{4},
\]
direct expansion gives
\[
 \det(V-tV^{\mathsf T})=A(t-1)^2+t,
 \qquad
 \Delta_{P(p,q,r)}(t)=1+A(t+t^{-1}-2).
\]
Thus $pq+pr+qr=-1$ gives Alexander polynomial one, determinant one, and
signature zero: the determinant of $V+V^{\mathsf T}$ is $4A-1=-1$, so its
two eigenvalues have opposite signs.  The characterization by this equation
for three odd parameters is established mathematics
\cite[Section~2.3]{belousov-pretzel}.

In particular, for every odd $a\ge3$ let
\begin{equation}\label{eq:mon:family}
 p=-a,\qquad q=a+2,\qquad r=\frac{a^2+2a-1}{2}.
\end{equation}
All three parameters are odd and have absolute value at least three.  Also
$p+q=2$ and $pq=-a(a+2)$, whence
$pq+pr+qr=-a(a+2)+2r=-1$.  The input is a knot by
Lemma~\ref{lem:mon:parity}; its three genuine denominators establish
nontriviality by Proposition~\ref{prop:mon:three}.  Consequently its Seifert
genus is exactly one, although both the Alexander polynomial and the
signature give the same answers as for the unknot.  The first example is
$P(-3,5,7)$, with
$V=\left(\begin{smallmatrix}1&3\\2&6\end{smallmatrix}\right)$.
This is a useful regression family for the arithmetic branch. With the
implementation's positive-denominator normalization, its signed
Montesinos integer $D$ is $+1$; the value $pq+pr+qr=-1$ uses signed
pretzel parameters. Their absolute values agree, as required.  The classical
alternative $(-a,2a-1,2a+1)$ has the same Alexander-triviality property;
neither family is claimed here as a new class of knots.

\subsection{Closure-independent forbidden tangles}

An essential distinction is that a tangle may occur inside a much more general
knot diagram.  The appropriate local property is whether the tangle can occur
inside \emph{any} unknot, not whether one selected closure is knotted.

\begin{theorem}[Nogueira--Salgueiro local obstruction, with input checks]
\label{thm:mon:local}
Let
\[
 T=[e]+R(\beta_1/\alpha_1)+\cdots+R(\beta_r/\alpha_r),
 \qquad \alpha_i\ge2,\quad\gcd(\alpha_i,\beta_i)=1,
\]
be a genuine two-string tangle with no internal closed components and
$r\ge2$.  Then $T$ embeds in some unknot if and only if $r=2$ and
\begin{equation}\label{eq:mon:congruence}
 F:=\beta_1\alpha_2+\beta_2\alpha_1
 \equiv1\ \text{or}\ -1\pmod{\alpha_1\alpha_2}.
\end{equation}
Therefore $r\ge3$, or failure of \eqref{eq:mon:congruence} when $r=2$,
is a nontriviality certificate for every knot containing $T$.
\end{theorem}

\begin{proof}
For $e=0$ this is \cite[Theorem~4.9(a)]{nogueira-salgueiro}; integral twists
can be absorbed into a rational summand, changing $F$ by a multiple of
$Q=\alpha_1\alpha_2$ and leaving the criterion unchanged.
Here is the cover argument identifying the exact obstruction.
The cover of the tangle ball has base disk and exceptional orders
$\alpha_1,\ldots,\alpha_r$.  Its boundary is incompressible for $r\ge2$.
For example, its group has the relations $[h,x_i]=1$ and
$x_i^{\alpha_i}h^{\beta_i}=1$, without the last relation of
\eqref{eq:mon:sfs-group}.  In the quotient by $h$ it is the free product
$C_{\alpha_1}*\cdots*C_{\alpha_r}$, and the boundary base loop maps to
$(x_1\cdots x_r)^{-1}$, an element of infinite order by reduced normal form.
The fiber $h$ itself has infinite order, since the homomorphism to
$(\mathbb Q,+)$ given by $h\mapsto1$ and
$x_i\mapsto-\beta_i/\alpha_i$ respects the relations.  Thus the two
boundary generators are independent and the boundary torus group injects.
In particular, the tangle is essential: a disk separating its two strings
would lift to a compression disk of that boundary torus.

If an essential two-string tangle occurs in an unknot, its complementary
tangle is rational \cite[Proposition~4.4]{nogueira-salgueiro}.
A finite complementary slope with denominator at least two would add another
genuine exceptional fiber; Proposition~\ref{prop:mon:three} rules this out.
The infinite complementary slope is ruled out separately: its closure is the
denominator closure of the horizontal sum, whose determinant is
$\prod_i\alpha_i>1$.  This product also follows by expressing the
denominator closure as the connected sum of the denominator closures of the
rational pieces.  Hence the only possible case is $r=2$ with an integral
complement $[u]$.  Its signed homology determinant is
\[
 F+(e+u)Q.
\]
The resulting link is rational and is the unknot precisely when this integer
is $\varepsilon\in\{1,-1\}$, giving exactly
\eqref{eq:mon:congruence}.  Conversely a successful sign gives the explicit
integral unknotting complement $u=(\varepsilon-F)/Q-e$.
Since $Q\ge4$, at most one of the two signs works.
\end{proof}

The genuine-tangle hypothesis can either be checked by following strands or
by verifying that at most one denominator is even.  It cannot be dropped
merely because the syntax contains rational summands.  The arithmetic witness
for the two-summand rejection branch is simply the recomputed residue
$F\bmod Q$, which must be neither $1$ nor $Q-1$.  No factorization of $Q$ is
needed.  Satisfying the congruence gives \emph{no positive verdict for an
arbitrary surrounding knot}: it only says that some choice of complement
could unknot this tangle.

For example, the tangle $R(1/3)+R(1/4)$ has $F=7$, $Q=12$, so it is
forbidden in every unknot.  Its numerator and denominator determinants are
$7$ and $12$, with gcd one.  Thus the obstruction is stronger on this example
than the gcd-of-two-closures determinant obstruction: neither local homology
order alone supplies the universal obstruction.
Likewise the first two boxes of the family \eqref{eq:mon:family} already give
the forbidden tangle $R(-1/a)+R(1/(a+2))$: here $F=-2$ and
$Q=a(a+2)\ge15$.  This local witness remains valid if the entire third-box
region is replaced by any exterior tangle giving a one-component closure.

\subsection{A rigorous planar certificate for occurrence in a PD diagram}

We now specify the missing link between arithmetic data and a supplied knot
diagram.  A list of fractions, or an abstract four-edge cut, is not by itself
an occurrence certificate.

Let $H$ be the expanded connected shadow of the tangle constructed by the
declared rational-tangle operations.  It has four degree-one port vertices;
its other vertices are crossings with a cyclic order and an overpassing
opposite pair.  Let $G$ be the validated connected plane shadow of the input
knot.  A certificate maps the crossing darts of $H$ injectively to the crossing
darts of $G$.  It must satisfy all the following requirements:
\begin{enumerate}
\item Crossing images are distinct; each crossing map preserves its cyclic
      order and the overpassing pair.  The same global orientation convention
      is used throughout.
\item Every internal edge pairing of $H$ equals the corresponding edge pairing
      in $G$.  Every dart at an image crossing is accounted for, either by an
      internal edge or by one of the four ports.  There are no unrecorded
      attachments or ignored incident darts.
\item Each of the four port darts belongs to a distinct edge of $G$ whose
      other endpoint is outside the selected crossing set. None of these
      edges is an internal edge of the pattern. This deliberately strict
      condition is the initial implementation contract; allowing two cut
      points on one exterior edge is a possible extension, not an accepted
      certificate here.
\item In the plane ribbon graph $H$, all four ports lie on one facial boundary,
      in the declared boundary order.  This is checked by the face permutation
      of the rotation system after adding the four degree-one vertices.
\item Following opposite darts through all crossings of $H$ gives exactly two
      open strands and no closed strand component.
\end{enumerate}
The builder's rational fractions are recomputed from its operations.  Matching
only fractions, crossing counts, or abstract unembedded graphs does not check
these requirements.

\begin{lemma}[The common-face certificate yields a separating disk]
\label{lem:mon:disk}
Under the preceding conditions there is a disk in the projection sphere whose
interior contains precisely the certified tangle diagram and whose boundary
meets the knot in its four declared ports.  In particular, the input knot
contains the declared tangle, up to the boundary-preserving diagram isotopy.
\end{lemma}

\begin{proof}
Embed $H$ in the projection sphere using the certified dart map, placing a cut
point in each of the four distinct outgoing edges. Take a small regular
neighborhood of this connected
embedded graph.  The boundary component corresponding to the declared face
is incident to all four outgoing port segments.  Every component of the
remaining diagram attaches to $H$ at a port: otherwise it would be a separate
component of the connected shadow $G$.  Since no edge crosses $H$, all such
components lie in that one face.  The other boundary components of the regular
neighborhood therefore enclose no part of the remaining diagram.
Fill those other boundary components with their complementary disks.  This
produces a disk; adjust its remaining boundary near the four terminal port
segments so that it meets each once and no other diagram edge.  Its port
order is the certified one.  Thickening this planar disk slightly in the
projection direction gives the required tangle ball in $S^3$.
\end{proof}

\begin{proposition}[Verification and deterministic pattern search]
\label{prop:mon:matching}
For a fixed connected tangle pattern with $m$ crossings and an input plane
knot diagram with $n$ crossings, the occurrence certificate can be verified
in $O(n+m)$ word operations after parsing.  All its orientation-preserving
dart occurrences can be found in $O(nm)$ word operations, plus parsing and
output costs.  A conservative bit bound adds a factor $O(\log(n+m))$.
\end{proposition}

\begin{proof}
All edge, rotation, strand, injectivity and face checks are traversals of
finite permutations and graphs.  For the search, choose one crossing of the
connected pattern as root.  Its image and cyclic offset have exactly $2n$ potential
choices: offsets zero and two preserve the overpassing opposite pair.  An internal edge then forces the image of the next crossing and
the image of its incoming dart; preservation of cyclic order forces its
remaining three darts.  Propagate this rule through the connected pattern,
rejecting inconsistent revisits, collisions, and overpair violations.
There are no further choices after the root choice.  Each candidate takes
$O(m)$ work, including the common-face and port checks.  The global validation
of $G$ is performed once. Direct indexed arrays give the stated word bound;
indices have $O(\log(n+m))$ bits.
\end{proof}

For a supplied already compiled template, this is the matching cost.
The public implementation additionally validates and relabels the target PD;
its sorting step gives $O(n\log(n+2)+nm)$ word operations with bounded labels.
Source reading, arithmetic, and template compilation have their own $B$- and
$m$-dependent cost and must be added if a template was not already compiled.
In particular, an arbitrarily long source with zero coefficients is not
processed in time bounded by $m$ alone.

The pattern expansion cost is explicit here.  If the chosen literal
continued-fraction pattern has more than $n$ crossings, it cannot have such
an injective occurrence in the supplied PD and may be discarded before
expansion.  This does not affect the independent compressed-input theorem,
which never expands twist exponents.  Searching a library of patterns with
total expanded size $M$ costs $O(nM)$ by the same argument.  Efficient
recognition of all Montesinos tangles hidden behind arbitrary diagram
isotopies is a different, unsolved implementation obligation in this package.

In particular, a seven-crossing template for $R(1/3)+R(1/4)$ gives an
$O(n\log(n+2))$-time public rejection stage (with a linear matching core) on every diagram containing
an admissible occurrence with four distinct outgoing cut edges.  This is an explicit family with unrestricted
exterior diagram size and complexity.  It is not a statement that every
nontrivial knot has such a tangle, or that finding it after arbitrary isotopies
has linear complexity.

The useful integration direction is one-sided: a verified occurrence plus
Theorem~\ref{thm:mon:local} gives a rigorous \texttt{KNOTTED} verdict.
Failure to find a pattern, failure of a resource budget, or arithmetic that
allows an unknotting closure causes the normal pipeline to continue.

\subsection{Bounded-description discovery without a supplied sidecar}

The preceding matching procedure also gives an explicit discovery bound when
the local tangle has a short literal description.  This removes the need for
a user-supplied occurrence certificate on a precisely delimited class.  It
does not require a conjecture about the descriptions of all nontrivial knots.

For integers $n\ge1$ and $h\ge2$, let $\mathcal S(n,h)$ consist of all
ordered lists of at least two nonempty continued-fraction blocks
\[
 s=\bigl([a_{1,1},\ldots,a_{1,\ell_1}],\ldots,
          [a_{b,1},\ldots,a_{b,\ell_b}]\bigr),\qquad b\ge2,
\]
such that
\[
 \sum_j\ell_j\le h,\qquad -n\le a_{j,k}\le n,
 \qquad \sum_{j,k}|a_{j,k}|\le n.
\]
Its specified diagram is obtained by the fixed rational-tangle compiler and
horizontal composition, with initial integer term $e=0$.  No diagram
isotopies are performed during matching.  The compiler expands an integer
$a$ into $|a|$ crossings and implements tangle rotation and reflection
without adding crossings.  Thus its literal pattern has at most $n$
crossings.  Zero coefficients and repeated descriptions are allowed; they
only make the following upper bound more generous.

Call a source \emph{admissible and forbidden} when every original final
fraction is finite and nonintegral with denominator greater than one, its
compiled shadow is connected, all four ports pass the common-face test, its
strand tracing gives exactly two open strands and no closed component, and
Theorem~\ref{thm:mon:local} rules out every unknotting complement. Occurrences
must have four distinct external cut edges, as above. This matches the strict
source grammar of the delivered local compiler: initial $e=0$, at least two
nonintegral rational summands, and no discarded integral summands. Integer
parts within a nonintegral slope do not affect its congruence obstruction.
The graph being matched is always the original literal compiled graph.

\begin{proposition}[Deterministic discovery from bounded source length]
\label{prop:mon:bounded-discovery}
Given a validated $n$-crossing plane knot diagram $G$ and an integer $h\ge2$,
there is a deterministic one-sided search with running time
\[
 (n+2)^{O(h)}
\]
in bit operations which returns \texttt{KNOTTED} if and only if $G$ contains
an orientation-preserving dart occurrence, satisfying
Lemma~\ref{lem:mon:disk}, of a pattern compiled from some admissible and
forbidden source in $\mathcal S(n,h)$.  When it returns \texttt{KNOTTED},
it also returns the source and the verified crossing-dart map.  Otherwise it
returns \texttt{UNKNOWN}.  In particular, for
$h\le C\log(n+2)$ with fixed $C$, this restricted witness-discovery problem
has running time
\[
 \exp\bigl(O((\log(n+2))^2)\bigr)
       =(n+2)^{O(\log(n+2))}.
\]
\end{proposition}

\begin{proof}
A source with total coefficient count $\ell$ is specified by a sequence of
$\ell$ coefficients and by selecting separators among the $\ell-1$ gaps.
There are $(2n+1)^\ell$ sequences and at most $2^{\ell-1}$ separator sets;
requiring at least two blocks and imposing the expansion budget can only
decrease their number.  Consequently
\[
 |\mathcal S(n,h)|
 \le \sum_{\ell=1}^{h}(2n+1)^\ell 2^{\ell-1}
 \le \bigl(2(2n+1)\bigr)^{h+1}.
\]
Enumerate these finite sequences and separators, and immediately discard any
source exceeding the expansion budget.  Evaluate every fraction by its
integer matrix product, reject integral final summands, and apply the exact
component and local obstruction tests.  The source bit length is
$B=O(h\log(n+2))$; all these arithmetic operations cost $B^{O(1)}$ by the
previous estimates.  Unsupported final infinite fractions are discarded.

Compile each remaining source, checking connectedness, its ports, and strand
components.  Standard explicit tangle composition and index permutations
are polynomial in the source and expanded lengths.  Even a compiler which
revisits the current graph at every continued-fraction operation takes
$O(hn)$ word operations, absorbed by
$O(n^2+B^2)$ up to logarithmic index costs.  Discard any inadmissible result.
For an admissible forbidden pattern with $m\le n$ crossings, run the
root-dart propagation of Proposition~\ref{prop:mon:matching}, costing
$O(nm)\le O(n^2)$ word operations.  Including integer arithmetic and bit
indices, with no fixed implementation expansion cap, the work per enumerated source is bounded by
\[
 O\bigl(n^2\log(n+2)+B^{C_0}\bigr)
\]
for some fixed constant $C_0$.  Multiplication by the preceding source count
is $(n+2)^{O(h)}$; polynomial factors in $h$ are absorbed since $h\ge2$.

Any returned map passes the separating-disk verification and the exact local
obstruction theorem, so its verdict is sound for the surrounding knot.
Conversely, every source in the stated finite class is enumerated.  If its
literal pattern has such an occurrence, the matching procedure tries the
correct root image and cyclic offset and then reconstructs the entire map.
It therefore finds that occurrence.  Substituting
$h\le C\log(n+2)$ gives the displayed quasi-polynomial bound.
\end{proof}

This is a theorem about discovering local certificates of bounded literal
description length in the supplied diagram.  Its completeness class is
specified by the source grammar and the checked plane embedding, rather than
by knot type alone.  A diagram requiring Reidemeister moves before such a
pattern becomes visible need not be rejected by the search.  No bound on
$h$ for arbitrary nontrivial knot diagrams is proved or assumed.  The
seven-crossing fixed-pattern matcher is an immediately usable implementation
of the underlying search primitive; the exhaustive source enumerator in
Proposition~\ref{prop:mon:bounded-discovery} is a separate algorithmic
construction, and should be distinguished from that implementation in
benchmark claims. The code defaults to a 100,000-crossing expansion cap;
that resource policy is not part of the abstract completeness theorem.
A fixed finite pattern library and a bounded exhaustive
enumerator may both precede the existing exact fallback recognizer.

